% ── sections/01_introduction.tex ──────────────────────────────────────────────

\section{Introduction}

Cancer arises through the accumulation of somatic mutations in a small
subset of genes---termed \textit{cancer driver genes}---that confer selective
growth advantages to affected cells. Distinguishing these drivers from the vast
background of passenger mutations is essential for understanding tumourigenesis,
stratifying patients, and selecting targeted therapies. Genome-wide sequencing
projects such as The Cancer Genome Atlas (TCGA) have catalogued thousands of
mutations across dozens of cancer types, yet the challenge of prioritising
functionally relevant alterations at scale remains largely unsolved
\cite{emogi2021,emgnn2023}.

Graph neural networks (GNNs) operating on protein--protein interaction (PPI)
networks offer a natural framework for this problem: genes with similar functions
cluster in network neighbourhoods, providing biologically meaningful relational
structure \cite{gcn2017}. Early methods such as EMOGI \cite{emogi2021} encode
each gene with multi-omics features and propagate information through a single
PPI network. EMGNN \cite{emgnn2023} extends this to a multilayer setting by
constructing a meta-graph that aggregates evidence from multiple per-network GNN
encoders. However, prior work reports a fundamental challenge for these
approaches: \textbf{PPI-based cancer driver prediction can be affected by
heterophily}, with neighbours of cancer drivers often being non-drivers
\cite{sgcd2024}. This suggests that standard message-passing GNNs, which
implicitly assume homophily, may dilute the cancer signal through neighbourhood
aggregation.
DISHyper \cite{dishyper2024} and DGHNN \cite{dghnn2025} partially address this
by incorporating functional annotations via hypergraph structures, while
deepCDG \cite{deepcdg2025} uses cross-omics attention, but none systematically
quantify or exploit the heterophily structure inherent in PPI topology.

This heterophily challenge manifests in three concrete open questions. First,
which architectural components---residual connections, normalisation strategies,
label smoothing, learning-rate schedules---genuinely improve performance on
biologically sparse, full-batch graph tasks, and which are harmful? Second, to
what extent does multi-network integration benefit prediction, and can we
disentangle the contribution of additional data from architectural improvements?
Third, can pathway-level interpretability connect model predictions to
established cancer hallmarks, providing biological validation beyond metric
improvement?

Here we address these questions through a systematic study built around the
EMGNN framework. Our core technical insight is that \textbf{heterophily-aware
gating}---a learned fusion of low-pass (neighbourhood-smoothed) and high-pass
(self-vs-neighbourhood deviation) filtered signals---effectively exploits the
heterophilous structure of PPI networks without requiring architectural redesign.
We implement this within a two-stage architecture (Fig.~\ref{fig:architecture}):
$K$ PPI networks are independently processed by a shared $L$-layer GCN encoder;
node embeddings are scaled by learnable per-network importance weights
$w_k = \mathrm{softmax}(\theta_k)$; and weighted embeddings are aggregated into
a meta-graph, processed by a second GCN, and classified. We systematically
evaluate this framework through five methodologies: (1) faithful benchmark
reproduction across three GNN backbones and six PPI databases, (2) controlled
ablation identifying BatchNorm as harmful in full-batch graph settings
($-$4.2\% AUPR), (3) multi-network extension with gain decomposition
(data +5.4\% vs.\ architecture +1.0\%), (4) biological validation via Integrated
Gradients and gene set enrichment (28 significant Hallmark pathways), and
(5) ablation of nine advanced GNN techniques, where heterophily-aware gating
achieves the strongest single-technique gain (AUPR~=~0.824~$\pm$~0.004, +2.8\%).
